\documentclass[11pt]{article}

%============引用宏包 和 自定义命令========================%
\usepackage{CJK}
\usepackage{graphicx}
\usepackage{latexsym,bm}
\usepackage{amsmath}
\usepackage{xcolor}
\usepackage{fancyhdr}
\pagestyle{fancy}

%============ 重定义字体、字号命令 =======================%
\newcommand{\hei}{\CJKfamily{hei}}

\begin{document}
\begin{CJK}{UTF8}{gkai}
%------------------------文章标题---------------------------------------------------------
\rhead{\tt L.J.SHOU  \includegraphics[scale=0.2]{ICT.eps} }  %页眉
\title{ \textbf{Assignment 6 of Algorithm Design and Analysis}}
\author{L.J.SHOU  寿林钧\\{\slshape No}. \bf 201228013229133}
\date{Dec.24 2012}
\maketitle % 生成题头
\noindent
{\bf \large  2 Problem Reduction}\\[2mm]
Let M be an n × n matrix with each entry equal to either 0 or 1. Let mij denote
the entry in row i and column j. A diagonal entry is one of the form mii for
some i.\\
Swapping rows i and j of the matrix M denotes the following action: we
swap the values mik and mjk for k = 1, 2, ..., n. Swapping two columns is defined
analogously.\\
We say that M is rearrangeable if it is possible to swap some of the pairs of
rows and some of the pairs of columns (in any sequence) so that, after all the
swapping, all the diagonal entries of M are equal to 1.\\
\indent
(a) Give an example of a matrix M that is not rearrangeable, but for which
at least one entry in each row and each column is equal to 1.\\
\indent
(b) Give a polynomial-time algorithm that determines whether a matrix M
with 0-1 entries is rearrangeable.\\
解：\\
(a)
$\left(
  \begin{array}{ccc}
    0 & 1 & 0 \\
    0 & 1 & 0 \\
    1 & 0 & 1 \\
  \end{array}
\right)$
is not rearrangeable.\\
(b) 令\\
\begin{eqnarray*}
V=\{v_1,v_2,...,v_n,v_{11},v_{22},...,v_{nn}\}\cup\{s,t\}\\
E_0=\cup_{i=0}^n\{(s,v_i)\},\\
E_1=\cup_{i=0}^n\{(v_i,v_{kk})|m_{ik}=1\},\\
E_2=\cup_{i=0}^n\{(v_{ii},t)\},\\
E=E_0\cup E_1 \cup E_2,\\
c(e)=
\begin{cases}
	1;     &\forall e\in E_0\cup E_2,\\
	\infty;&\forall e\in E_1 \\
\end{cases}
\end{eqnarray*}

定义最大流问题 $P = (V, E, c, s, t)$. 如果问题 P 的最大流函数 $f$ 满足 $f (s, V ) = n$, 则矩
阵 $M$ 是可重排的,则是不可重排的。因此只需要按照上面的方法,
将问题转化成最大流问题,再用求解最大流问题的算法求解即可。\\[2mm]

{\bf \large  3 Unique Cut}\\

Let$ G = (V, E)$ be a directed graph, with source $s\in V$ , sink$ t\in V$ , and
nonnegative edge capacities ${c_e}$. Give a polynomial-time algorithm to decide
whether G has a {\bf unique} minimum $s − t$ cut.\\
{\bf 解}:\\
应用 Ford-Fulkerson 算法计算最大流,
并得到最小割 $(S, T )$. 对每一条边 $e\in E(G)$ 且 $e = (u, v)$,$ u\in S$,$ v \in T$ . 将$ (u, v)$ 的容量增加一个单位得新的网络流 $G_e$,
如果对所有新的网络流的最大流有 $f_{G_e} (s, V ) > f_G (s, V )$, 则唯一,
否则不唯一。\\



\end{CJK}
\end{document}

